% !TeX root = ../PIFB2.tex
\section{Gauge Covariance, Empirical Symmetry and Scientific Agreement}\label{sec:gauge_invariance_consensus}

Scientific statements must distinguish a change in the system from a change in its description. Gauge covariance addresses the latter distinction. Consensus addresses whether different participants make compatible statements about records. An empirical symmetry addresses whether specified transformations of a physical preparation and its measurements preserve predictive relations. These questions can interact, but they are not equivalent.

\subsection{Coordinate consistency does not select a physical symmetry}

For a bimeasurable invertible change of representation $g$, pushing both a law and its reference forward preserves their relative entropy,
\begin{equation}
D_{\mathrm{KL}}(g_\#q\|g_\#p)=D_{\mathrm{KL}}(q\|p).
\label{eq:science_coordinate_consistency}
\end{equation}
The same consistency requires links, record kernels, noise and priors to transform with their declared types. It holds for one agent before any agreement with another agent occurs. Consensus therefore does not create this invariance. Changing a coordinate frame while holding a coordinate-dependent noise covariance or likelihood fixed generally changes the model as well.

\begin{remark}[Covariance and selection]
\label{rem:kretschmann}
An invertible re-expression of one model cannot select a physical law over its re-expression, because it changes no predictive probabilities. A claim that variational dynamics select a physical symmetry must instead specify distinct candidate models or states and a gauge-invariant criterion distinguishing their predictions or stability.
\end{remark}

The Lorentz construction illustrates the distinction precisely. Once a complex rank-two communication sector is chosen, its determinant-preserving complex congruences have the proper, orthochronous Lorentz group as their image. This theorem supplies a representation and its invariant. It does not make every selective interaction a Lorentz transformation, select that channel from an arbitrary cognitive space, or establish that observed signals propagate along its null cone. Likewise, embedding a familiar compact group in a sufficiently large general linear group would show representational availability, not derive the interaction law of a physical gauge field.

Empirical symmetry needs an operational relation. If $T$ acts on latent preparations and $R$ is a bimeasurable bijection on records, a proposed covariant record mechanism must satisfy
\begin{equation}
K(RA\mid Ty)=K(A\mid y)
\label{eq:operational_symmetry_kernel}
\end{equation}
for every measurable record event $A$ and every preparation $y$ in its declared domain, together with compatible laws for preparation and evolution. The relation becomes testable when $T$ and $R$ correspond to specified interventions and calibrated comparisons. If they merely rename every symbol, Eq.~\eqref{eq:operational_symmetry_kernel} is a consistency identity. If the physically realizable interventions obey it approximately, their measured residual and regime of validity are part of the physical claim.

\subsection{Relational invariants of the communication network}

Independent links carry content that pure coordinate frames do not. For a closed directed path at one fixed context, form the composable operator
\begin{equation}
\mathcal L_\gamma=L_{i_0i_1}L_{i_1i_2}\cdots L_{i_{r-1}i_0}:E_{i_0}\longrightarrow E_{i_0}.
\label{eq:science_loop_operator}
\end{equation}
Endpoint re-expressions cancel at intermediate agents, giving $\mathcal L_\gamma'=g_{i_0}\mathcal L_\gamma g_{i_0}^{-1}$. Its rank, characteristic polynomial and traces of powers are therefore invariant. This calculation does not require each link to be invertible. For singular links the product is a loop operator, not an element of a transport group; for invertible group-valued links it can be called network holonomy. Neither construction alone is curvature of a connection along paths in the base.

These invariants suggest targeted relational observations: compare responses to different composable communication routes, controlling the endpoint law and output channel. A nontrivial internal loop operator is observable only to the extent that those records distinguish its action; a law invariant under that action can hide it. This is why zero transported disagreement does not force identity holonomy and why changing a graph loop need not change a particular loss. The required claim is about a specified observable response, not a universal relation between an invariant's magnitude and performance.

\subsection{Agreement, evidence and model comparison}

The normalized population law separates shared evidence from shared conclusions. A record appears once in its generative factorization even when several agents can access it. Selective communication can promote agreement on a low-dimensional public statistic while retaining different private states. A community that agrees is therefore not, on that fact alone, calibrated or predictively successful. Priors, source preferences and restrictions on record access can sustain agreement about an incorrect model.

The evidential comparison is predictive. A candidate should specify record probabilities before the relevant outcome, compare them with those of an alternative under the same access constraints, and assess interventions that were not used to fit the retained description. In the monograph's mathematical examples, the frozen-source contrast, interface-rank changes and removal of retained memory isolate different mechanisms. Their theoretical predictions are not empirical results about nature until a corresponding system and record protocol are supplied.

This also limits the cognitive reading of scientific change. A passive change of representational frame is scientifically empty. A change in latent variables, conditional independence assumptions, accessible records or predictive kernels can alter what a model explains. Describing both changes as a new gauge would conceal the substantive one. The general framework permits a rigorous comparison only after those altered model components have been specified.

\subsection{A falsifiable participatory program}

The claim that every possible observation can be redescribed as a consequence of cognitive architecture has no discriminating prediction in that unrestricted form. The present program instead asks for mechanisms with a restricted domain: an actual normalized model, independently specified observations, a retained variable and a forecast about its response or composition. Conditional mathematical results identify which hypothesis can fail. A parent-link square may not exist, a retained law may require memory, a proposed geometry may be degenerate, or its causal characteristics may disagree with its determinant cone. These are informative failures of specified constructions.

The underlying contextual base remains part of the model, and observer access remains mediated by records and inference. Neither fact proves that an observer-independent reality is absent or that scientific success is exhausted by agreement. The participatory perspective directs attention to how stable shared descriptions are formed and tested. Its physical content will be whatever quantitative relations survive that testing, with their assumptions and limits retained.

